% БҮЛЭГ 3 — СИСТЕМИЙН АРХИТЕКТУР БА ЗАГВАРЫН ЗОХИОН БАЙГУУЛАЛТ

\chapter{Системийн архитектур ба загварын зохион байгуулалт}
\label{Chapter3}

Энэ бүлэгт туршилтын архитектур, сургалтын дамжлага, интерфейсийн бүтэц,
үнэлгээний арга зүйг тайлбарлана.

\textbf{Чухал тэмдэглэл.} Энэ бүлэгт хоёр шатлалт каскад архитектурыг
танилцуулах боловч каскад алдааны дамжуулалтын улмаас $78.61\%$ гүйцэтгэл
гаргасан тул \textbf{эцсийн шийдэл болгон нэг шатлалт Flat MN-BERT-ийг сонгосон}
(харьцуулалтыг бүлэг~\ref{Chapter5}–\ref{Chapter6}-д үзнэ үү).

%-------------------------------------------------------------------------------
\section{Системийн туршилтын архитектур}

Системийн зохион байгуулалтыг илүү тодорхой болгохын тулд эхлээд системийн
хэрэглээний тохиолдлын (Use Case) болон үйл ажиллагааны (Activity) логикийг
тодорхойлов.

\subsection{Хэрэглээний тохиолдлын диаграмм (Use Case Diagram)}

Зураг \ref{fig:use-case}-т харуулсанчлан систем нь гурван үндсэн оролцогчтой:
Өсвөр насны хэрэглэгч, Модератор (зохицуулагч) болон Системийн администратор.
Хэрэглэгч сэтгэгдэл оруулах бол, Модератор нь системийн ангилсан үр дүнг хянаж,
хортой агуулгыг шүүх, бүтээлч шүүмжлэлийг зөвшөөрөх үйлдлийг гүйцэтгэнэ.

\begin{figure}[H]
\centering
\begin{tikzpicture}[
  usecase/.style={ellipse, draw, minimum width=5.0cm, minimum height=1.0cm,
                  fill=blue!8, draw=blue!50, line width=0.8pt,
                  align=center, font=\small},
  actor/.style={circle, draw, minimum size=0.7cm,
                fill=orange!14, draw=orange!65, line width=0.8pt},
  arr/.style={draw, ->, >=Stealth, thick, line width=0.7pt},
]

  % ── SYSTEM BOUNDARY ──────────────────────────────────────────────────────
  \draw[rounded corners=6pt, draw=gray!55, fill=gray!4, line width=1pt]
    (-3.2, -4.6) rectangle (3.2, 3.6);
  \node[font=\small\bfseries, text=gray!65] at (0, 3.25)
    {NLP Ангиллын Систем};

  % ── USE CASES ────────────────────────────────────────────────────────────
  \node[usecase] (uc1) at (0,  2.2) {Сэтгэгдэл оруулах};
  \node[usecase] (uc2) at (0,  0.8) {Ангиллын дүн үзэх};
  \node[usecase] (uc3) at (0, -0.7) {Хортой агуулга шүүх};
  \node[usecase] (uc4) at (0, -2.2) {Бүтээлч шүүмжлэл зөвшөөрөх};
  \node[usecase] (uc5) at (0, -3.8) {Загвар сургах / шинэчлэх};

  % ── LEFT ACTORS ──────────────────────────────────────────────────────────
  \node[actor] (user)  at (-6.0,  2.2) {};
  \node[font=\small, align=center] at (-6.0, 1.35)
    {Өсвөр насны\\Хэрэглэгч};

  \node[actor] (admin) at (-6.0, -3.8) {};
  \node[font=\small, align=center] at (-6.0, -4.65)
    {Системийн\\Администратор};

  % ── RIGHT ACTOR ──────────────────────────────────────────────────────────
  \node[actor] (mod)   at ( 6.0, -0.7) {};
  \node[font=\small, align=center] at ( 6.0, -1.55)
    {Модератор};

  % ── CONNECTIONS ──────────────────────────────────────────────────────────
  \draw[arr] (user)  -- (uc1.west);
  \draw[arr] (admin) -- (uc5.west);
  \draw[arr] (mod)   -- (uc2.east);
  \draw[arr] (mod)   -- (uc3.east);
  \draw[arr] (mod)   -- (uc4.east);

\end{tikzpicture}
\caption{Системийн хэрэглээний тохиолдлын диаграмм: Хэрэглэгч сэтгэгдэл оруулж,
  Модератор ангиллын үр дүнг хянан шүүж, Администратор загварыг удирдана.}
\label{fig:use-case}
\end{figure}


\subsection{Үйл ажиллагааны диаграмм (Activity Diagram)}
\label{sec:activity-diagram}

Системийн логик урсгалыг Зураг \ref{fig:activity-diagram}-т харуулав. Энэхүү
диаграмм нь хоёр шатлалт ангиллын шийдвэр гаргах үйл явцыг нарийвчлан харуулж
байна.

\begin{figure}[H]
    \centering
    \resizebox{\linewidth}{!}{%
    \begin{tikzpicture}[
        node distance=1cm and 1.5cm,
        start/.style={circle, draw, fill=black, inner sep=3pt},
        stop/.style={circle, draw, double, fill=black!50, inner sep=2.5pt},
        procstep/.style={rectangle, draw, rounded corners, minimum width=4cm, minimum height=1cm, text centered, fill=green!5, drop shadow, align=center},
        decision/.style={diamond, draw, aspect=2, minimum width=3cm, text centered, fill=yellow!10, font=\small, drop shadow, inner sep=0pt, align=center},
        arrow/.style={thick, ->, >=stealth}
    ]
        \node (init) [start] {};
        \node (input) [procstep, below=of init] {Текст хүлээн авах};
        \node (prep) [procstep, below=of input] {Урьдчилсан боловсруулалт};
        \node (s1) [decision, below=1.5cm of prep] {Сөрөг? (Stage 1)};
        
        \node (allow) [procstep, left=2.5cm of s1] {Зөвшөөрөх (Allow)};
        
        \node (s2) [decision, below=1.5cm of s1] {Хортой? (Stage 2)};
        \node (block) [procstep, right=2.5cm of s2] {Блок хийх (Block)};
        
        \node (end) [stop, below=of s2] {};

        \draw [arrow] (init) -- (input);
        \draw [arrow] (input) -- (prep);
        \draw [arrow] (prep) -- (s1);
        \draw [arrow] (s1) -- node[anchor=east, xshift=-0.1cm] {Тийм} (s2);
        \draw [arrow] (s1) -- node[anchor=south] {Үгүй} (allow);
        \draw [arrow] (s2.west) -| node[anchor=south, xshift=1.5cm] {Үгүй} (allow);
        \draw [arrow] (s2) -- node[anchor=south] {Тийм} (block);
        
        \draw [arrow] (allow) |- (end);
        \draw [arrow] (block) |- (end);
    \end{tikzpicture}}
    \caption{Системийн шүүлтүүрийн логикийн үйл ажиллагааны диаграмм}
    \label{fig:activity-diagram}
\end{figure}

\subsection{Давхаргын архитектур}

Системийн архитектур нь өгөгдлийн боловсруулалт, загварын давхарга, хэрэглэгчийн интерфейс гэсэн гурван давхаргаас бүрдэнэ.

\begin{figure}[H]
  \centering
  \begin{tikzpicture}[
      block/.style={rectangle, draw, fill=blue!10, text width=10em,
                    text centered, rounded corners, minimum height=3em, align=center},
      line/.style={draw, -latex}
    ]
    \node [block] (input) {Хам текст \\ (найман эх сурвалж)};
    \node [block, below of=input, node distance=3cm] (preprocess)
          {Урьдчилсан боловсруулалт \\ (Цэвэрлэгээ, Токенчлол)};
    \node [block, below of=preprocess, node distance=3cm] (stage1)
          {Stage~1: Сентимент ангилал \\ (Эерэг / Саармаг / Сөрөг)};
    \node [block, below of=stage1, node distance=3cm] (stage2)
          {Stage~2: Нарийвчилсан ангилал \\ (Хортой сөрөг vs Бүтээлч шүүмжлэл)};
    \node [block, below of=stage2, node distance=3cm] (output)
          {Үр дүн \\ (Gradio интерфейс)};
    \path [line] (input) -- (preprocess);
    \path [line] (preprocess) -- (stage1);
    \path [line] (stage1) -- node [right, xshift=0.3cm]
          {\small Сөрөг} (stage2);
    \path [line] (stage2) -- (output);
  \end{tikzpicture}
  \caption{Системийн ерөнхий архитектурын схем}
  \label{fig:system-arch}
\end{figure}

Зураг~\ref{fig:pipeline}-д дамжлагыг харуулав.

\begin{figure}[H]
\centering
\begin{tikzpicture}[
  node distance=0.5cm,
  box/.style={draw, rounded corners, minimum width=4.8cm, minimum height=0.75cm,
              text centered, font=\small, fill=blue!8},
  wide/.style={draw, rounded corners, minimum width=4.8cm, minimum height=0.75cm,
               text centered, font=\small, fill=green!12},
  arrow/.style={->, >=Stealth, thick}
]
\node[box]  (raw)   at (0,0)    {Хам текст (news.mn / Facebook / gogo.mn)};
\node[box]  (filt)  at (0,-1.3) {Кирилл шүүлт ($\geq$85\% кирилл)};
\node[box]  (clean) at (0,-2.6) {Цэвэрлэгээ (URL, HTML, emoji → орлуулагч)};
\node[box]  (tok)   at (0,-3.9) {SentencePiece токенчлол (MN-BERT)};
\node[box]  (enc)   at (0,-5.2) {\texttt{[CLS]}+токен дараалал → 12-давхарга энкодер};
\node[wide] (cls)   at (0,-6.5) {\texttt{[CLS]} вектор ($\mathbb{R}^{768}$) → Шугаман + Softmax};
\node[wide] (out)   at (0,-7.8) {Ангиллын шошго (Эерэг / Саармаг / Хортой сөрөг / Бүтээлч шүүмжлэл)};
\draw[arrow] (raw)   -- (filt);
\draw[arrow] (filt)  -- (clean);
\draw[arrow] (clean) -- (tok);
\draw[arrow] (tok)   -- (enc);
\draw[arrow] (enc)   -- (cls);
\draw[arrow] (cls)   -- (out);
\end{tikzpicture}
\caption{MN-BERT сургалтын дамжлага (pipeline) (хам текстээс ангиллын шошго хүртэл)}
\label{fig:pipeline}
\end{figure}

Stage~1: Эерэг/Саармаг/Сөрөг ангилал; зөвхөн Сөрөг текст Stage~2-т
Хортой сөрөг/Бүтээлч шүүмжлэл ялгалтанд орно (Зураг~\ref{fig:two-stage-arch}).

%-------------------------------------------------------------------------------
\section{Хоёр шатлалт ангиллын архитектур}

Хоёр шатлалт архитектур: Эерэг/Саармаг/Сөрөг ялгалт (Stage~1) болон Хортой
сөрөг/Бүтээлч шүүмжлэл нарийвчлал (Stage~2). Лексик шинжилгээ хангалтгүй тул
контекст мэдрэмжтэй загвар шаардагдах бөгөөд хоёр шат бүрийг тусгай даалгавартаа төвлөрүүлнэ.

\subsection{Stage~1 --- Сентимент ангилал}

Stage~1 загварын зорилго нь оролтын текстийг Эерэг, Саармаг, Сөрөг гэсэн
гурван ангилалд хуваах юм. Энэ шатанд дараах загваруудыг хэрэгжүүлж,
харьцуулна:

\begin{itemize}
  \item \textbf{MN-BERT:} нарийн тохируулга (fine-tune) хийж \texttt{[CLS]} гаралтыг softmax руу дамжуулна.
  \item \textbf{LR:} TF-IDF дээр L2 регуляризацитай.
  \item \textbf{NB:} TF-IDF дээр Multinomial Naive Bayes.
  \item \textbf{SVM:} TF-IDF дээр Шугаман цөм (Linear kernel).
\end{itemize}

Stage~1-ийн гол шаардлага: Сөрөг ангиллын сэргээлт (recall) өндөр байлгах (алдаагүй тохиолдолд Stage~2-т дамжуулна).

\subsection{Stage~2 --- Хортой сөрөг vs Бүтээлч шүүмжлэл ялгалт}

Stage~1-ээс Сөрөг гэж ангилагдсан текстүүд Stage~2 руу дамжина. Энэ шатанд
хоёр ангилалт (binary classification) хийнэ:

\begin{itemize}
  \item \textbf{Хортой сөрөг:} Доромжлол, хувь хүнд чиглэсэн довтолгоо --- устгах шаардлагатай.
  \item \textbf{Бүтээлч шүүмжлэл:} Аргументтай, нотолгоонд суурилсан сөрөг үнэлгээ --- хамгаалах шаардлагатай.
\end{itemize}

Stage~2-т MN-BERT болон суурь загваруудыг (LR, NB, SVM) харьцуулна.
Ялгалтын шалгуур: (1)~\textbf{чиглэл} (хувь хүн vs. асуудал);
(2)~\textbf{зорилго} (гомдоох vs. сайжруулах) --- MN-BERT-ийн анхаарлын механизм (self-attention) тэр ялгааг барина.

\subsection{Хоёр шатлалт загварын давхаргын архитектур}

Хоёр шатлалт загварын давхаргын бүтцийг Зураг~\ref{fig:two-stage-arch}-д харуулав.
Stage~1 ба Stage~2 нь ижил MN-BERT encoder, харин өөр өөр classifier head-тай.

\begin{figure}[H]
\centering
\scalebox{0.76}{%
\begin{tikzpicture}[
  layer/.style={draw, rounded corners=4pt, minimum width=7.5cm,
                minimum height=1.40cm, text width=7.0cm,
                align=center, font=\small, line width=0.7pt},
  hlayer/.style={draw, rounded corners=4pt, minimum width=3.5cm,
                 minimum height=1.40cm, text width=3.2cm,
                 align=center, font=\small, line width=0.7pt},
  arr/.style={->, >=Stealth, thick, line width=0.8pt},
]

% ── Shared MN-BERT backbone ───────────────────────────────────────────────
\node[layer, fill=gray!12, draw=gray!50] (inp) at (0, 0)
  {\textbf{Оролт токенчлол}\\[3pt]
   {\scriptsize [CLS] $t_1 \ldots t_n$ [SEP] \enspace·\enspace max\_length\,=\,128}};

\node[layer, fill=yellow!16, draw=yellow!65!black] (emb) at (0,-2.60)
  {\textbf{Эмбэддинг давхарга (768-хэмжээст)}\\[3pt]
   {\scriptsize Token Emb + Position Emb + Segment Emb}};

\node[layer, fill=blue!12, draw=blue!50] (enc) at (0,-5.20)
  {\textbf{БЕРТ Энкодер --- 12 Transformer блок}\\[3pt]
   {\scriptsize Multi-Head Self-Attention (12 толгой) + LayerNorm + Feed-Forward}};

\node[layer, fill=blue!22, draw=blue!60] (cls) at (0,-7.80)
  {\textbf{[CLS] Pooled вектор --- 768-хэмжээст}};

\node[layer, fill=orange!12, draw=orange!55] (drop) at (0,-10.40)
  {\textbf{Dropout}\enspace(p\,=\,0.1)};

% ── Stage 1 head (left) ───────────────────────────────────────────────────
\node[hlayer, fill=orange!20, draw=orange!60] (lin1) at (-2.6,-13.60)
  {\textbf{Шугаман}\\[3pt]{\scriptsize 768\,$\to$\,3}};

\node[hlayer, fill=green!16, draw=green!55] (out1) at (-2.6,-16.20)
  {\textbf{Softmax (Stage~1)}\\[3pt]
   {\scriptsize Эерэг · Саармаг · Сөрөг}};

% ── Stage 2 head (right) ──────────────────────────────────────────────────
\node[hlayer, fill=teal!14, draw=teal!55] (lin2) at (2.6,-13.60)
  {\textbf{Шугаман}\\[3pt]{\scriptsize 768\,$\to$\,2}};

\node[hlayer, fill=red!14, draw=red!50] (out2) at (2.6,-16.20)
  {\textbf{Softmax (Stage~2)}\\[3pt]
   {\scriptsize Хортой сөрөг · Бүтээлч шүүмжлэл}};

% ── Shared backbone arrows ────────────────────────────────────────────────
\draw[arr] (inp)  -- (emb);
\draw[arr] (emb)  -- (enc);
\draw[arr] (enc)  -- (cls);
\draw[arr] (cls)  -- (drop);

% Fork to two heads (0.80cm gap before branching)
\draw[arr] (drop.south) -- ++(0,-0.80) coordinate (fork)
           -- node[above, font=\scriptsize, text=green!55!black] {\textit{Stage~1}}
           (-2.6,-11.90) -- (lin1.north);
\draw[arr] (fork)
           -- node[above, font=\scriptsize, text=teal!55!black] {\textit{Stage~2}}
           (2.6,-11.90) -- (lin2.north);

\draw[arr] (lin1) -- (out1);
\draw[arr] (lin2) -- (out2);

\end{tikzpicture}%
}
\caption{Хоёр шатлалт MN-BERT загварын давхаргын архитектур. Хоёр загвар тус бүр
  ижил 12-блокт БЕРТ энкодерыг (768-хэмжээст hidden state) ашигладаг боловч
  тус тусдаа нарийн тохируулга (fine-tuning) хийгдсэн.
  Stage~1 нь 3-ангиллын (Шугаман 768$\to$3),
  Stage~2 нь 2-ангиллын (Шугаман 768$\to$2) ангиллын толгойтой (classifier head).}
\label{fig:two-stage-arch}
\end{figure}

%-------------------------------------------------------------------------------
\section{MN-BERT загварын нарийн тохируулга}

MN-BERT нь Монгол кирилл корпус дээр урьдчилан сургагдсан (pre-train) БЕРТ загвар.
Эцсийн гиперпараметрийг Хүснэгт~\ref{tab:hyperparams-mnbert} (бүлэг~\ref{Chapter5})-д үзнэ.

\subsection{Нарийн тохируулга стратеги}

Stage~1 нь 3-ангиллын cross-entropy, Stage~2 нь binary cross-entropy алдагдлын (loss) функцээр тус тусдаа нарийн тохируулга хийгдэнэ; хоёулаа ижил урьдчилан сургагдсан жинг эхлэлийн цэг болгоно. Early stopping: validation loss 2 epoch буурахгүй бол зогсооно. Class-weighted loss: цөөн ангилалд (minority class) урвуу давтамжийн жин тооцоолно.

%-------------------------------------------------------------------------------

\subsection{Загварын гиперпараметрүүд}
\label{sec:hyperparams-detail}

Гиперпараметр: LR $2\times10^{-5}$, багц (batch) 16, max\_length 128, AdamW (decay $0.01$), warmup 10\%. Эцсийн нэг шатлалт загварын тохиргоог Хүснэгт~\ref{tab:hyperparams-mnbert}-аас үзнэ үү.

\section{Суурь загварууд}

MN-BERT-ийн гүйцэтгэлийг дараах гурван суурь загвартай харьцуулна:

\begin{table}[H]
\centering
\caption{Суурь загваруудын тохиргоо}
\label{tab:baseline-config}
\begin{tabular}{@{}lll@{}}
\toprule
\textbf{Загвар} & \textbf{Шинж чанар} & \textbf{Тохиргоо} \\
\midrule
Logistic Regression & TF-IDF (unigram + bigram) & L2 регуляризаци, $C=1.0$ \\
Naive Bayes         & TF-IDF (unigram + bigram) & Multinomial NB \\
SVM                 & TF-IDF (unigram + bigram) & Шугаман цөм, $C=1.0$ \\
\bottomrule
\end{tabular}
\end{table}

Эдгээр суурь загварууд нь MN-BERT-ийн нэмүү үнэ цэнийг тодорхойлох жишиг сорил (benchmark) болно.

Тэнцвэргүй өгөгдлийн шийдлийг бүлэг~\ref{Chapter4} ба бүлэг~\ref{Chapter5}-д, Gradio хэрэглэгчийн интерфейсийн дэлгэрэнгүй бүтцийг бүлэг~\ref{Chapter5}-д тус тус авч үзнэ.

%-------------------------------------------------------------------------------

%-------------------------------------------------------------------------------
\section{Тэнцвэргүй өгөгдлийн шийдэл}
\label{sec:imbalance}

Датасетийн ангиллын тэнцвэргүй байдлыг хоёр стратегиар шийдвэрлэв.
Суурь загварт SMOTE-оор цөөн ангиллыг нийлэгжүүлж тэнцвэржүүлэв.
Эцсийн нэг шатлалт MN-BERT загварт Focal Loss ($\gamma{=}2.0$),
ангиллын жин $\alpha$, WeightedRandomSampler-ийн хослолыг ашиглав.
Энэхүү хослол хортой/бүтээлч ялгалтын цөөн ангиллын сэргээлтийг
мэдэгдэхүйц сайжруулснаар Focal Loss-ын онолын давуу тал батлагдав.


%-------------------------------------------------------------------------------
\section{Gradio хэрэглэгчийн интерфейс}
\label{sec:gradio}

Gradio фреймворк ашиглан хэрэглэгчийн интерфейс бүрдүүлэв. Хэрэглэгч
текст оруулахад систем урьдчилсан боловсруулалт хийж, нарийн тохируулсан MN-BERT
загварт дамжуулан Эерэг/Саармаг/Хортой сөрөг/Бүтээлч шүүмжлэл гэсэн
4 ангиллын нэгийг итгэлцлийн оноотой хамт буцаана.

Инференсийн урсгал: $\text{Текст} \to \text{Урьдчилсан боловсруулалт} \to
\text{MN-BERT токенчлол} \to \text{Нэг шатлалт ангилал} \to \text{Хэрэглэгч}$.
Нэг шатлалт загвар ($83.26\%$) ба Хоёр шатлалт ($78.61\%$) хоёулаа нэгдсэн
интерфейст байдаг бөгөөд хэрэглэгч сонгох боломжтой.

\section{Үнэлгээний арга зүй}

Гүйцэтгэлийг Нарийвчлал (Accuracy), Оновчлол (Precision), Сэргээлт (Recall), $F_1$, Confusion Matrix-аар хэмжинэ.
Тэнцвэргүй өгөгдлийн нөхцөлд \textbf{Макро дундаж $F_1$ (Macro-average $F_1$)}-ийг гол үзүүлэлт болгох бөгөөд
статистик найдвартай байдлыг McNemar тест болон bootstrap CI-аар нотолно.

\subsection{Харьцуулалтын арга зүй}

Загвар бүрийг ижил туршилтын өгөгдөл (test set) дээр үнэлнэ: ижил train/val/test хуваалт, ижил
урьдчилсан боловсруулалт (MN-BERT: tokenizer; суурь: TF-IDF); туршилтын өгөгдөл сургалтанд ашиглагдаагүй;
Макро дундаж $F_1$-score гол үзүүлэлт.

%-------------------------------------------------------------------------------

\subsection{Ангиллын үзүүлэлтүүд}
\label{sec:eval-metrics}

Үнэлгээнд Нарийвчлал, Макро-Оновчлол (macro-Precision), Макро-Сэргээлт (macro-Recall), Макро-F$_1$
(\textit{macro}: ангилал тус бүрийн жинлэгдсэн дундаж) болон
Confusion Matrix ашиглав. Тэнцвэргүй өгөгдөлд Макро F$_1$ нь
нарийвчлалаас (accuracy) илүү мэдрэмжтэй хэмжүүр тул эцсийн шалгуур болно.

\section{Технологийн стек}

Системийг хэрэгжүүлэхэд ашиглагдах гол технологиудыг хүснэгт~\ref{tab:tech-stack}-д нэгтгэв.

\begin{table}[H]
\centering
\caption{Системийн технологийн стек}
\label{tab:tech-stack}
\begin{tabular}{@{}ll@{}}
\toprule
\textbf{Зориулалт} & \textbf{Технологи} \\
\midrule
Програмчлалын хэл       & Python 3.10+ \\
Загварын framework       & PyTorch, Hugging Face Transformers \\
Урьдчилсан загвар       & MN-BERT (Hugging Face Hub) \\
Суурь загварууд          & scikit-learn \cite{pedregosa2011} \\
Тэнцвэргүй өгөгдөл      & imbalanced-learn (SMOTE) \\
Хэрэглэгчийн интерфейс  & Gradio \\
Өгөгдлийн боловсруулалт  & pandas, NumPy \\
Дүрслэл                  & matplotlib, seaborn \\
Хувилбар хяналт          & Git, GitHub \\
\bottomrule
\end{tabular}
\end{table}


%-------------------------------------------------------------------------------
\section{Бүлгийн дүгнэлт}

Энэ бүлэгт архитектур, нарийн тохируулгын стратеги, суурь загварууд, тэнцвэржүүлэлт, Gradio интерфейс, үнэлгээний арга зүйг тодорхойлсон.
